\section{Proposed methodology}
\subsection{Search strategy}
We will conduct a structured scoping review to map concepts, methods and
evidence gaps across heterogeneous vision applications. This design fits
the three descriptive subquestions and the need to compare assessment
prerequisites; it does not estimate a single pooled treatment effect.
IEEE Xplore, ACM Digital
Library and CVF proceedings will be the core sources; Scopus will be used
if institutional access is available. Citation tracing and targeted arXiv
searches will supplement the search, with preprints marked separately.
The proposed scope is English-language publications from 2017 to the final
search date, with earlier foundational sources included through citation
tracing. The final search date will be recorded when the search is completed.

The initial query will combine (synthetic data OR synthetic images OR
simulated data) AND (computer vision OR image classification OR object
detection OR semantic segmentation) AND (quality OR evaluation OR realism
OR fidelity OR diversity OR coverage OR bias OR sim-to-real).
We will adapt syntax to each database and run companion searches for
requirements, named metrics and assessment prerequisites. Before screening,
we will pilot queries against known relevant sources and record exact queries,
filters, dates, result counts and protocol changes.

\subsection{Selection and quality appraisal}
Studies will be eligible if they evaluate synthetic vision training data,
identify requirements for its use, or introduce assessment methods relevant
to the research questions. We will exclude unrelated modalities, duplicate
versions and sources without sufficient detail to interpret their assessment.
Domain-general metric papers will be classified as methodological background.

After deduplication by DOI and title, both authors will independently screen
titles and abstracts, followed by full texts. We will discuss disagreements,
retain the individual and agreed decisions, and record full-text exclusion
reasons. Appraisal will consider task relevance, independence of real test
data, annotation quality, baseline comparability, subgroup reporting,
uncertainty and reproducibility. These judgements will qualify the synthesis
rather than be combined into an unexplained numerical score.

\subsection{Extraction and synthesis}
One author will extract each study and the other will verify it. Records
will include task, domain, generation method, real and synthetic datasets,
data splits, training regime and source locations for claims.
For RQ1, we will extract requirements concerning realism, similarity,
coverage, labels, bias and downstream utility. For RQ2, we will record metric
definitions, assumptions, implementation settings and limitations.
For RQ3, we will identify required real images, annotations, priors, trained
networks, human expertise and computational resources.

We will compare studies within compatible tasks and protocols, using a
narrative synthesis and a requirement--metric--resource matrix. Scores from
different datasets, representations or training budgets will not be pooled
without a comparability rationale. Missing evidence, conflicting results and
the distinction between study findings and our interpretation will be explicit.
